% Abhakseite „Das kann ich“ – Zone nur im Gesamt (\mitzone)
%% TODO Ich-kann-Sätze: Platzhalter wie in den Aufgabendateien
\begin{abhakseite}
\ifdefined\mitzone
\abhakgruppe{Kennst du schon}
\abhak{1}{Natürliche Zahlen addieren, subtrahieren, multiplizieren, dividieren im Kopf (Einmaleins)}
\abhak{2}{Zahlenstrahl}
\abhak{3}{Dezimalzahlen und Brüche addieren, subtrahieren, multiplizieren (Kommazahlen addieren, Bruch mal ganze Zahl)}
\abhak{4}{Punkt vor Strich und Klammern mit natürlichen Zahlen}
\abhak{5}{Wert eines Terms mit Platzhalter berechnen (Vorzahl mal x plus Zahl, für einen gegebenen Einsetzwert)}
\abhak{6}{Punkt vor Strich und Klammern mit natürlichen Zahlen – Fehler finden}
\abhak{7}{Punkt vor Strich und Klammern mit natürlichen Zahlen – selbst rechnen}
\fi
\abhakgruppe{Einheit 1 · Negative Zahlen kennen und ordnen}
\abhak{8}{Ordnen}
\abhak{9}{Mitte zweier Zahlen}
\abhak{10}{runden}
\abhak{11}{Punkte in vier Quadranten}
\abhak{12}{Ordnen – Fehler finden, Begründen, Darstellungswechsel, Anwendung}
\abhakgruppe{Einheit 2 · Addieren und Subtrahieren}
\abhak{13}{Vorzeichen oder Rechenzeichen?}
\abhak{14}{Zeichen zusammenfassen}
\abhak{15}{Addieren/Subtrahieren}
\abhak{16}{Beträge: gleiche Vorzeichen addieren, verschiedene subtrahieren}
\abhak{17}{Dezimalzahlen und Brüche}
\abhak{18}{Umkehrung: Rechnung zu Ergebnis finden}
\abhak{19}{Addieren/Subtrahieren – Fehler finden, Begründen, Darstellungswechsel, Anwendung}
\abhakgruppe{Einheit 3 · Multiplizieren und Dividieren}
\abhak{20}{Multiplizieren/Dividieren}
\abhak{21}{Vorzeichenregel als Aussage prüfen}
\abhak{22}{Multiplizieren/Dividieren – Fehler finden, Begründen, Anwendung}
\abhakgruppe{Einheit 4 · Terme und Sachaufgaben}
\abhak{23}{Terme und Sachaufgaben}
\abhak{24}{Rechenvorteile (Tauschen, geschickt zusammenfassen)}
\abhak{25}{Terme und Sachaufgaben – Fehler finden, Begründen, Anwendung}
\end{abhakseite}
